\subsection{The \acrlong{lartpc}} \label{sec:LARTPC}
The \gls{lartpc} technology is based on the use of liquid argon as both the target and detection medium for neutrino interactions. It was proposed by Carlo Rubbia in 1977~\cite{rubbia1977liquid} becoming a leading technology for neutrino detection due to its excellent spatial resolution, calorimetric capabilities, and scalability. In a \gls{lartpc}, when a neutrino interacts with an argon nucleus, it produces charged particles that ionize the argon atoms along their path. The liberated electrons drift under the influence of an electric field towards a readout plane, where they are collected and recorded as signals. The readout plane is typically segmented into multiple layers (induction and collection planes) that allow for a 2D reconstruction of the event. By combining the information from these planes with timing information, a full 3D reconstruction of the neutrino interaction can be achieved. Fulfilling the fundamental principle of a \gls{tpc}. Additionally, liquid argon produces scintillation light when ionized, which can be detected by a \gls{pds}.

\begin{figure}[ht!]
    \centering
    \includegraphics[width=.8\linewidth,clip,trim=0cm 4cm 0cm 2cm]{Images/Schematics/ORIGINAL_TPC.pdf}
  \caption[width=1\linewidth]{Schematic of the original liquid argon time projection chamber (LArTPC) showing the readout plane~\cite{rubbia1977liquid}.\label{fig:LARTPC}}
\end{figure}

%\gls{dune}'s \gls{fd} the coordinate system is designed such that the collection plane has its readout wires or strips oriented along the beamline direction ($z$ coordinate). The complementary \gls{tpc} coordinate $y$ is reconstructed from a redundant matching of two induction planes. The arrival times and charges of all signal are recorded, allowing for a projected 2D-reconstruction of the interaction $< t, y, z >$. In addition, the scintillation light can be associated to the \gls{tpc} information for non-beam triggering, improved calorimetry, background rejection and drift coordinate $x$ reconstruction. For low-energy events, efficient charge collection and combined light detection and matching are crucial for a reliable 3D-reconstruction and background suppression.